\section{Address relations and stable families}
\label{sec:core-relations}

Let $\Sigma_m=\{1,\ldots,m\}^{\mathbb N}$, and let
$\pi_p:\Sigma_m\to K_p$ be the coding map of a contractive similarity
system. We use the following standard implication of the connectedness
criterion to specify connected families.

\begin{proposition}[Connected families from prescribed contacts]
\label{prop:prescribed-contacts}
Suppose a parameter family satisfies finitely many coding equalities
$\pi_p(i\omega)=\pi_p(j\eta)$ whose first-symbol pairs contain the
edges of a connected graph on $\{1,\ldots,m\}$. Then every attractor
in the family is connected.
\end{proposition}
\begin{proof}
The prescribed equalities give edges in the first-piece intersection
graph. To see the connectedness criterion directly, start with the
convex hull of $K_p$ and iterate the union operator, as in
Proposition~\ref{prop:binary-contact}. At each stage the images are
connected and have a connected intersection graph, since they contain
the corresponding pieces of $K_p$. The nested intersection is $K_p$.
\end{proof}

This separates enforcing connectedness from understanding additional
contacts. The complex-tree constructions in
\cite{Espigule2013,Espigule2018,Espigule2019a,Espigule2019b}
motivate treating the latter as a parameter-space problem.

\begin{proposition}[Regularity of address equations]
\label{prop:address-regularity}
For a fixed plane-orientation type, $\pi_p(\omega)$ is real analytic
in the normalized coefficients. It is holomorphic in the
direct--direct class. If $\omega$ is eventually periodic, its value
is a rational function of the coefficients and their conjugates;
in the direct--direct class it is a holomorphic rational function.
Consequently eventually periodic contact equations define relatively
real-algebraic sets, and in the direct--direct class complex algebraic
sets, after denominators are cleared. A contact equation gives a
smooth real surface in the four-dimensional chart at points where
its real Jacobian has rank two. No such rank assertion is automatic.
\end{proposition}

\begin{proof}
Write each map as $t_j+L_j$, where $t_j$ is its translation and
$L_j$ is its real-linear similarity part. The coding series is
\[
\pi_p(\omega)=t_{\omega_1}
+\sum_{n\ge1}L_{\omega_1}\cdots L_{\omega_n}t_{\omega_{n+1}}.
\]
On a compact contraction domain its terms have a geometric uniform
bound. For opposite similarities, complexify the coefficients and
their conjugates as independent variables. The same series converges
normally on a sufficiently small complex neighbourhood where all
these variables have modulus less than one. Restriction to the real
parameter domain proves real analyticity. In the direct--direct
case no conjugate variable occurs, proving holomorphic dependence.

For an address $uv^\infty$, its value is
$S_u(\Fix S_v)$. A direct affine map $c+\nu z$ has fixed point
$c/(1-\nu)$; an opposite map $c+\nu\overline z$ has fixed point
$(c+\nu\overline c)/(1-|\nu|^2)$. Word compositions have polynomial
coefficients in the original coefficients and their conjugates, and
their fixed-point denominators do not vanish under strict
contraction. This proves the rational and algebraic assertions.
The final assertion is the real implicit function theorem.
\end{proof}

For nonperiodic addresses the equations are generally analytic, not
polynomial. Nor does a contact equation by itself identify a piece
of $\partial\cM$: it can lie inside $\cM$, be a forced identity on
a connected slice, or have singular and lower-rank points.

% Optional short section: exact contact seeds, with their core location proved.
\subsection{Explicit eventually periodic contacts in the direct class}
\label{sec:restored-periodic-contacts}

Consider $f_-(z)=-1+az$ and $f_+(z)=1+bz$, with
$0<|a|,|b|<1$. Their fixed points are
$p_-=-1/(1-a)$ and $p_+=1/(1-b)$.

\begin{proposition}[Cross-contact algebraic families]
\label{prop:restored-periodic-contacts}
For every integer $k\ge1$,
\begin{equation}\label{eq:restored-periodic-criterion}
 f_-^k(p_+)=f_+^k(p_-)
 \quad\Longleftrightarrow\quad a^k+b^k=1.
\end{equation}
These parameter sets are smooth complex curves within the strict
contraction bidisk and every attractor on them is connected.
For $k\ge2$ they lie in the universal connected core
$|a|^2+|b|^2\ge1$. For $k>2$ this inequality is strict, so all such
parameters are in the interior of the full connectedness locus.
For $k=2$, equality holds only for real $a,b$.
\end{proposition}

\begin{proof}
Geometric summation and subtraction give
\[
 f_-^k(p_+)-f_+^k(p_-)
 =-\frac{(a+b-2)(a^k+b^k-1)}{(1-a)(1-b)}.
\]
The denominator is nonzero, and $a+b\ne2$ in the strict contraction
bidisk. The two points are coded by $-^k+^{\infty}$ and
$+^k-^{\infty}$, respectively, so their equality gives first-piece
contact. The gradient $(ka^{k-1},kb^{k-1})$ is nonzero, proving
smoothness of the complex curve.

If $a^k+b^k=1$, the triangle inequality gives
$|a|^k+|b|^k\ge1$. For $k\ge2$ and both moduli strictly between
zero and one,
\[
 |a|^2+|b|^2\ge |a|^k+|b|^k\ge1,
\]
with strictness in the first inequality for $k>2$.
For $k=2$, equality throughout requires $a^2$ and $b^2$ to be
nonnegative real numbers, hence $a,b$ to be real.
\end{proof}

In the atlas coordinates define
\[
 S_k(\alpha,w)=w^ke^{ik\alpha}+(1-w)^ke^{-ik\alpha}.
\]
Where $S_k\ne0$, the same contact curves have branches
\[
 \rho=2|S_k|^{1/k},\qquad
 \gamma=\frac{\arg S_k+2\pi\ell}{k},
 \quad \ell=0,\ldots,k-1,
\]
retaining only points with $\rho>2\max(w,1-w)$ and imposing the
phase-lattice identifications. These are exact connected families
with explicit contacts. For $k>2$ they cannot be frontier sheets;
their role is to provide algebraic examples and contact seeds
inside the connected region.


\subsection{Mirror symmetry and algebraic contact curves}

\begin{proposition}[Mirror-paired contacts]\label{prop:mirror-contact}
In the direct--direct chart restrict to $b=\overline a$, with
$0<|a|<1$, and let $\overline\omega$ denote the address obtained
by exchanging the two symbols of $\omega$. Then
\[
\pi_a(\overline\omega)=-\overline{\pi_a(\omega)}.
\]
The contact $\pi_a(\omega)=\pi_a(\overline\omega)$ is therefore
equivalent to the single real equation
\begin{equation}\label{eq:mirror-contact}
\Re\pi_a(\omega)=0.
\end{equation}
Every parameter satisfying this equation belongs to $\cM$.
If $\omega$ is eventually periodic, equation~\eqref{eq:mirror-contact}
is real algebraic after a nonvanishing denominator is cleared.
At a point where the gradient of its left-hand side in
$(\Re a,\Im a)$ is nonzero, its zero set is a smooth real-analytic
curve; in the eventually periodic case it is a regular branch of
a real-algebraic curve.
\end{proposition}

\begin{proof}
The reflection $R(z)=-\overline z$ satisfies $RS_1R=S_2$ and
$RS_2R=S_1$. Applying these identities to finite words and then
passing to the coding limit gives the displayed symmetry.
A point is fixed by $R$ exactly when its real part is zero, proving
\eqref{eq:mirror-contact}. The two paired addresses have opposite
first symbols, so their equality puts a point in both first-level
pieces and implies connectedness. The algebraicity assertion follows
from Proposition~\ref{prop:address-regularity}; the gradient condition
and the real-analytic implicit function theorem give the local curve.
\end{proof}

This is a precise mechanism for smooth algebraic pieces in
mirror-symmetric contact families. It does not make every such
piece a connectedness-frontier branch: contact curves may lie in
the interior of the ambient connectedness locus. For several
simultaneous address equations on the two-real-dimensional mirror
slice, one must check their effective number of independent
constraints. A single regular scalar equation, or a system of
constant rank one in a neighbourhood, gives a curve. Rank one at
one point of an arbitrary system alone does not suffice. Rank two
can give isolated solutions, while rank zero requires a separate
singularity analysis.

In the historical common-translation mirror chart
$T_1(z)=1+az$, $T_2(z)=1+\overline a z$, ordinary conjugation exchanges
the maps and the corresponding paired contact equation is
$\Im\pi_a(\omega)=0$. For nonreal $a$, canonical normalization
turns that reflection into $R(z)=-\overline z$ and preserves the
contact condition. At real $a$ the two common-translation maps
coincide, so this historical chart degenerates to a singleton;
the canonical chart should not be identified with it there.

\begin{definition}[Forced equivalences and structural stability]
For a nonempty parameter family $X$, define
\[
\mathscr R_X=\bigcap_{p\in X}\Ker\pi_p,
\qquad
\mathcal S_X=\{p\in X:\Ker\pi_p=\mathscr R_X\}.
\]
Thus $\mathscr R_X$ consists of address identifications forced
throughout the family, and $\mathcal S_X$ consists of parameters
with no additional identifications. We call this structural
stability relative to $X$; it does not assert dynamical stability
or openness in the ambient parameter space.
\end{definition}

\begin{proposition}[Topology on the stable set]
\label{prop:kernel-quotient}
For every $p\in\mathcal S_X$, the coding map induces a homeomorphism
$\Sigma_m/\mathscr R_X\to K_p$. In particular all attractors in
$\mathcal S_X$ have the same topology. Neither nonemptiness nor
openness of $\mathcal S_X$ follows from its definition.
\end{proposition}
\begin{proof}
Each coding kernel is a closed equivalence relation on compact
$\Sigma_m$, so their intersection is one also. Its quotient is
compact Hausdorff. At a parameter in $\mathcal S_X$, the induced
continuous map to $K_p$ is bijective, hence a homeomorphism.
\end{proof}

